\documentclass{situationist-preprint}

\title{Ethics as First Philosophy and the Problem of the Third}
\msid{PHE-13-b}
\msfamily{PHE-13}
\msdiscipline{Phenomenology}
\mstype{Consequence}
\mskeywords{Levinas; first philosophy; the third party; justice; infinite demand; obligation}

\begin{document}
\maketitle

\begin{abstract}
Levinas claims that ethics is first philosophy and that the third party introduces justice. Critics see a dilemma: if the dyadic relation to the face is first, justice, which compares the incomparable, betrays it; if the third is always already present, the dyad is an abstraction and ethics is not first. We derive from PHE-13-a that the dilemma is false. Ethics is first in the order of normative sources: every obligation is a selection, made under measure, from claims that bind independently of it. The third is first in the order of occurrence: no claim is ever met outside comparison. Since the two priorities concern different orders, they are consistent, and justice neither betrays nor replaces the demand.
\end{abstract}

\section{The dilemma}

Levinas holds that the relation to the other, in which I am called into question by a face, precedes and conditions knowledge and ontology, so that ethics, not metaphysics, is first philosophy \citep{Levinas1961,Levinas1984}. He also holds that the third party is present from the start in the eyes of the other, and that the third introduces comparison, measure and justice \citep[V.3]{Levinas1974}. Readers have found these claims in tension. If the relation to the single face is first and unconditional, justice, which treats faces as comparable cases, appears as a betrayal of it; if the third is never absent, the pure dyad is never encountered, and the priority of ethics looks like the priority of a fiction \citep{Derrida1997,Critchley1992}.

PHE-13-a provides what is needed to state the claims precisely. The demand of a face is infinite and binding: every discharge leaves a binding claim unanswered (Theorem~3.2). Obligation is the domain of fault and is jointly satisfiable (Definition~2.6 and Assumption~2.7). Not every binding claim is obligatory (Theorem~3.4), and some discharge meets every obligation (Corollary~3.5). Section~4 of PHE-13-a describes the passage from the demand to obligation as a selection, made by the measure that the third introduces, of a jointly satisfiable set of binding claims.

\section{Two orders of priority}

\begin{definition}\label{d:source}
A normative status $N$ is \emph{prior in the order of sources} to a status $M$ if whatever has $M$ has it in virtue of having $N$ together with some further condition, while having $N$ does not depend on having $M$.
\end{definition}

\begin{definition}\label{d:occurrence}
A condition $C$ is \emph{prior in the order of occurrence} to a practice if every exercise of the practice takes place under $C$.
\end{definition}

\begin{proposition}\label{p:ethics}
Binding is prior to obligation in the order of sources.
\end{proposition}

\begin{proof}
On the account of Section~4 of PHE-13-a, a claim is obligatory if and only if it binds and is selected by measure into a jointly satisfiable set. So every obligation has its status in virtue of binding together with selection. Conversely, by Assumption~2.5 of PHE-13-a, a claim binds because of its source in the face, whether or not it is selected; Theorem~3.4 gives claims that bind and are not obligatory. So binding does not depend on obligation.
\end{proof}

\begin{proposition}\label{p:third}
The third party is prior to every obligation in the order of occurrence.
\end{proposition}

\begin{proof}
Obligation requires joint satisfiability (Assumption~2.7 of PHE-13-a), which is a relation among several claims weighed against one another and against the agent's capacities. Weighing is comparison, and comparison is what the third introduces. Since, as Levinas holds, the third is never absent, every obligation arises under comparison, and every answer an agent gives to a claim is given under it.
\end{proof}

\begin{theorem}\label{t:consistent}
Ethics is first philosophy in the order of sources, and the third is first in the order of occurrence; the two priorities are consistent.
\end{theorem}

\begin{proof}
The first conjunct is Proposition~\ref{p:ethics}: the normative force from which every obligation derives is that of the face, which is not derived from measure, capacities or principles. The second is Proposition~\ref{p:third}. Definition~\ref{d:source} concerns what a status depends on, Definition~\ref{d:occurrence} the conditions under which a practice is exercised. A status can be underived and yet never be acted on outside a given condition, so the two priorities cannot conflict.
\end{proof}

Theorem~\ref{t:consistent} answers both horns. Justice does not betray the face, because what justice selects has its force from the face and what it does not select continues to bind (Corollary~3.5 of PHE-13-a). The dyad is an abstraction in the order of occurrence, but abstraction is how a source is isolated: no one encounters a force apart from the bodies it acts on, and its priority in explanation is not thereby a fiction. Levinas's own formula, that justice is required by responsibility and must always be corrected by it, follows: required, because obligation exists only under measure; corrected, because the binding claims that a measure leaves out are standing grounds for revising it.

\section{Objections and replies}

\begin{objection}
If every answer is given under comparison, the infinite demand has no practical role; it is idle.
\end{objection}

\begin{reply}
Its role is critical. The unselected claims bind (Theorem~3.4 of PHE-13-a), and a binding claim is a reason, even when failing to answer it is no fault. Such reasons bear on the choice of measure itself, which is why a binding claim that is not an obligation today may become one under a better measure. A source that is never acted on outside comparison can still be what comparisons are answerable to.
\end{reply}

\begin{objection}
First philosophy in Levinas's sense is meant to precede ontology, not merely obligation. Theorem~\ref{t:consistent} concerns only the order among normative statuses.
\end{objection}

\begin{reply}
The present result is restricted, deliberately, to what PHE-13-a establishes. It shows that the priority of the face over justice is coherent; whether the ethical relation also precedes the understanding of being is a further thesis, which the framework of PHE-13-a neither asserts nor denies.
\end{reply}

\section{Conclusion}

The problem of the third is a problem only if the priority of ethics is read as a priority of occurrence. Read as a priority of sources, it is consistent with the omnipresence of the third and explains why justice is both demanded by the face and perpetually answerable to it.

\begin{thebibliography}{9}\small
\bibitem[Critchley(1992)]{Critchley1992} Critchley, S. (1992). \emph{The Ethics of Deconstruction: Derrida and Levinas}. Blackwell.
\bibitem[Derrida(1997)]{Derrida1997} Derrida, J. (1997). \emph{Adieu \`a Emmanuel L\'evinas}. Galil\'ee.
\bibitem[Levinas(1961)]{Levinas1961} Levinas, E. (1961). \emph{Totalit\'e et infini: Essai sur l'ext\'eriorit\'e}. Martinus Nijhoff.
\bibitem[Levinas(1974)]{Levinas1974} Levinas, E. (1974). \emph{Autrement qu'\^etre ou au-del\`a de l'essence}. Martinus Nijhoff.
\bibitem[Levinas(1984)]{Levinas1984} Levinas, E. (1984). Ethics as first philosophy. Trans. S. Hand and M. Temple. In S. Hand (Ed.), \emph{The Levinas Reader} (1989). Blackwell.
\end{thebibliography}

\end{document}
